\apendice{Anexo de sostenibilización curricular}

\section{Introducción}
El presente anexo recoge una reflexión personal sobre cómo los principios de sostenibilidad se han integrado y abordado durante el desarrollo de este Trabajo Fin de Grado (TFG). Más allá de los aspectos puramente técnicos, el proyecto "Laboratorio Virtual de Programación y Robótica" ha ofrecido la oportunidad de explorar y aplicar una visión de la ingeniería informática como una disciplina al servicio del desarrollo sostenible. Se profundiza en las competencias de sostenibilidad adquiridas y su aplicación práctica, demostrando que las soluciones tecnológicas pueden ser catalizadores para un futuro más equitativo y responsable, en línea con las directrices de la CRUE \cite{CRUE2012_Directrices}
\section{Contribución del proyecto a los Objetivos de Desarrollo Sostenible (ODS)}
El proyecto se alinea de manera fundamental con varios de los Objetivos de Desarrollo Sostenible de la ONU, especialmente en lo que respecta al acceso equitativo a la educación y la reducción de las desigualdades.
\subsection{ODS 4: Educación de calidad}
El \textbf{ODS 4: Educación de Calidad} busca garantizar una educación inclusiva, equitativa y de calidad, así como promover oportunidades de aprendizaje durante toda la vida para todos. El laboratorio virtual contribuye directamente a este objetivo al derribar barreras geográficas y económicas que limitan el acceso a la enseñanza de la programación y la robótica. La dependencia de laboratorios físicos bien equipados y kits de robótica específicos representa un obstáculo considerable para centros educativos con recursos limitados o en zonas remotas. Al ofrecer una alternativa de simulación en Unity, se democratiza el acceso a la educación STEAM (Ciencia, Tecnología, Ingeniería, Arte y Matemáticas), permitiendo a niños y adolescentes experimentar con la programación de robots sin necesidad de adquirir costosos materiales o infraestructura. Esto promueve la inclusión de un mayor número de estudiantes en disciplinas de alto valor añadido, fomentando el desarrollo del pensamiento computacional y las habilidades digitales cruciales para el siglo XXI, asegurando que el origen socioeconómico no sea un factor determinante en el acceso a la calidad educativa.
\subsection{ODS 10: Reducción de las desigualdades}
De forma inherente, el proyecto aborda el \textbf{ODS 10: Reducción de las desigualdades}, el cual aspira a reducir la desigualdad dentro y entre los países, facilitando la inclusión social, económica y política de todas las personas. El laboratorio virtual disminuye la brecha de acceso a la tecnología y al conocimiento especializado que a menudo existe en contextos de vulnerabilidad. Al priorizar el desarrollo para dispositivos Android tipo tablet, se garantiza que la plataforma sea accesible en una gama más amplia de hardware de bajo coste, extendiendo así la oportunidad de aprender programación y robótica a segmentos de la población que tradicionalmente quedarían excluidos. La gratuidad de las herramientas empleadas para su desarrollo (Unity Personal, Csharp, UBlockly de código abierto, Git/GitHub) reduce drásticamente el coste de implementación y replicación de esta solución, permitiendo que organizaciones educativas, comunidades y familias puedan acceder a recursos didácticos de vanguardia sin que la inversión económica sea un impedimento. Este enfoque minimiza las disparidades en las oportunidades de aprendizaje, contribuyendo a una sociedad más equitativa e inclusiva.
\section{Competencias de sostenibilidad adquiridas y aplicadas}
Durante el desarrollo de este TFG, he tenido la oportunidad de aplicar y fortalecer diversas competencias clave relacionadas con la sostenibilidad, trascendiendo el ámbito técnico y proyectándose hacia un impacto social y medioambiental.
\subsection{Pensamiento crítico y sistémico}
El desafío de crear un laboratorio virtual que sea accesible y educativo ha requerido un profundo pensamiento crítico y sistémico. Comprender las interconexiones entre la falta de recursos, las desigualdades educativas y la solución tecnológica implicaba analizar cómo cada componente del diseño (desde la elección de Unity y UBlockly hasta la simplificación del modelo de robot LEGO WeDo 2.0) impactaría en la accesibilidad global y en la viabilidad a largo plazo. Se ha procurado optimizar los algoritmos de simulación para un menor consumo de recursos, y el diseño de la interfaz para la usabilidad en pantallas pequeñas, evidenciando una perspectiva integral sobre cómo la tecnología se inserta en un contexto social y económico.
\subsection{Diseño orientado al impacto social y a largo plazo}
La reflexión sobre la sostenibilidad ha influido significativamente en las decisiones de diseño. En lugar de desarrollar una solución ad-hoc, se ha priorizado la creación de un sistema extensible y modular que pueda evolucionar. La elección de UBlockly y la arquitectura MVC responden a la necesidad de construir una base que pueda adaptarse a nuevos tipos de bloques, modelos de robot y plataformas (como futuras versiones web o de realidad virtual) sin tener que reinventar la rueda, prolongando la vida útil del software y maximizando su potencial impacto educativo a largo plazo. Esta visión a largo plazo minimiza el despilfarro de recursos de desarrollo y facilita la resiliencia del proyecto ante futuros cambios tecnológicos o necesidades pedagógicas.
\subsection{Gestión ética de la información y responsabilidad profesional}
La utilización de herramientas de código abierto como UBlockly y la contribución implícita a la comunidad de conocimiento libre forman parte de una gestión responsable y ética. Se ha prestado atención a la calidad y claridad del código, así como a la documentación técnica, para asegurar que el proyecto pueda ser mantenido y extendido por otros desarrolladores, fomentando la colaboración y el intercambio de conocimientos. Además, la conciencia sobre la importancia de la privacidad de los usuarios (aunque este prototipo no gestione datos sensibles, la mentalidad de diseñar para la seguridad futura es clave) y la integridad del proceso de desarrollo han sido principios rectores que han consolidado mi comprensión de la responsabilidad que implica el diseño de sistemas informáticos.
\section{Reflexión final}
El desarrollo de este TFG ha sido una experiencia transformadora que ha trascendido la mera aplicación de conocimientos técnicos en programación. Me ha permitido conectar la teoría de la ingeniería informática con desafíos sociales concretos y con los principios universales de la sostenibilidad. La convicción de que la tecnología bien diseñada, con un propósito claro y un enfoque ético, tiene un potencial inmenso para contribuir a un mundo más justo y equitativo se ha reforzado enormemente. Este proyecto es una muestra de cómo las competencias técnicas y de sostenibilidad se entrelazan para dar vida a soluciones innovadoras que abren nuevas oportunidades educativas y contribuyen al bienestar global. Este trabajo no es solo un ejercicio académico, sino el inicio de una vocación por construir un futuro donde la tecnología sirva como un verdadero motor de cambio sostenible.